\documentclass[a4paper,12pt]{article}
\usepackage{utils}

\title{ Géometrie 2 }
\date{\today}

\begin{document}
\maketitle
\tableofcontents
\clearpage

\section{Théorèmes classiques + barycentre}

\subsection{Les barycentres}

Dans cette section $\E$ est un espace affine réel et $\vec{\E}$ l'espace vectoriel associé.

\vspace{1em}
On rappelle la définition d'un espace affine.

On a $\funcdef{\theta}{\vec{\E} \times \E}{\E}{(\vec{v}, A)}{A + \vec{v}}$
une application telle que

\begin{itemize}
    \item $\forall A, B \in \E, \exists! \vec{v} \in \vec{\E}, B = A + \vec{v}$
    \item $\forall A \in \E, \forall \vec{u}, \vec{v} \in \vec{\E}, (A + \vec{u}) + \vec{v} = A + (\vec{u} + \vec{v})$
\end{itemize}

Une application $f: \E \to \F$ est affine s'il existe $\vec{f}: \vec{\E} \to \vec{\F}$ une application linéaire
et un point $O \in \E$ tels que $\forall \vec{v} \in \vec{\E}, f(O + \vec{v}) = f(O) + \vec{f}(\vec{v})$.

\vspace{1em}

\begin{definition}
    On considère les systèmes finis de points pondérés de $\E$

    $\mathcal{A} = ((A_1, \lambda_1), (A_2, \lambda_2), \ldots, (A_n, \lambda_n))$
    avec $(A_i, \lambda_i) \in \E \times \R$

    $\lambda = \sum_{i=1}^n \lambda_i$ est appelé \textbf{le poids total} de $\mathcal{A}$.

    On définit l'application
    $$\funcdef{\L}{\E}{\vec{\E}}{M}{\sum_{i=1}^n \lambda_i \overrightarrow{M A_i}}$$
    où $\overrightarrow{M A_i}$ est le vecteur tel que $A_i = M + \overrightarrow{M A_i}$

    Cette fonction s'appelle \textbf{la fonction de Leibniz}.
\end{definition}

\begin{theoreme}{}{}
    \begin{enumerate}
        \item $\lambda = 0 \implies \L$ est constant
        \item $\lambda \neq 0 \implies \L$ est bijective.
    \end{enumerate}
\end{theoreme}
\break
\begin{demonstration}
    \begin{enumerate}
        \item On calcule pour N et M dans $\E$
        \begin{align*}
            \L(M) - \L(N) &= \sum_{i=1}^n \lambda_i \overrightarrow{M A_i} - \sum_{i=1}^n \lambda_i \overrightarrow{N A_i} \\
            &= \sum_{i=1}^n \lambda_i \overrightarrow{MN} \\
            &= \lambda \overrightarrow{MN}
        \end{align*}

        Si $\lambda = 0$ alors $\L$ est constante.

        \item On a de même pour $\lambda \neq 0$

        \begin{align*}
            \L(M) = \L(N) &\iff \L(M) - \L(N) = \vec{0} \\
            &\iff \lambda \overrightarrow{MN} = \vec{0} \\
            &\iff \overrightarrow{MN} = \vec{0} \\
            &\iff M = N
        \end{align*}

        Il reste à montrer que $\L$ est surjective.

        Soit $\vec{v} \in \vec{\E}$
        Soit $O \in \E$ un point quelconque tel que $\L(O) = \vec{v_0}$
        On cherche $M \in \E$ tel que $\L(M) = \vec{v}$
        On a
        \begin{align*}
            \L(M) - \L(O) = \lambda \overrightarrow{MO} &\iff \vec{v} - \vec{v_0} = \lambda \overrightarrow{MO} \\
            &\iff \overrightarrow{MO} = \frac{1}{\lambda} (\vec{v} - \vec{v_0}) \\
            &\iff \overrightarrow{OM} = \frac{1}{\lambda} (\vec{v_0} - \vec{v}) \\
            &\iff M = O + \frac{1}{\lambda} (\vec{v_0} - \vec{v})
        \end{align*}
    \end{enumerate}
\end{demonstration}

\begin{definition}
    Le point $G$ tel que $\L(G) = \vec{0}$ est unique et s'appelle le \textbf{barycentre} de $\mathcal{A} = ((A_1, \lambda_1), \ldots, (A_n, \lambda_n))$
    On le note $G = bar(\mathcal{A})$
\end{definition}

Ainsi on a
\begin{align*}
    G = bar(\mathcal{A}) &\iff \sum_{i=1}^n \lambda_i \overrightarrow{G A_i} = \vec{0} \\
    &\iff \forall M \in \E, \sum_{i=1}^n \lambda_i (\overrightarrow{G M} + \overrightarrow{M A_i}) = \vec{0} \\
    &\iff \lambda \overrightarrow{G M} + \sum_{i=1}^n \lambda_i \overrightarrow{M A_i} = \vec{0} \\
    &\iff \lambda \overrightarrow{M G} = \sum_{i=1}^n \lambda_i \overrightarrow{M A_i} \\
    &\iff \overrightarrow{MG} = \frac{1}{\lambda} \sum_{i=1}^n \lambda_i \overrightarrow{MA_i}
\end{align*}

\begin{remarque}
    On peut écrire tous les points comme barycentres de $((A, \lambda), (B, \mu))$ pour $\lambda, \mu \in \R$.
\end{remarque}

\begin{proposition}{}{}
    Soit $\mathcal{A} = ((A_1, \lambda_1), \ldots, (A_n, \lambda_n)), \lambda \neq 0$
    et $G = bar(\mathcal{A})$
    \begin{enumerate}[label=\roman*.]
        \item Si on change l'ordre des points pondérés le barycentre ne change pas.
        \item Si on retire les points à poids nuls on ne change pas le barycentre.
        \item Si on multiplie tous les poids par $\alpha \neq 0$ alors le barycentre ne change pas et le poids total est $\alpha \lambda$.
        \item Pour $k \in \llbracket 1, n \rrbracket$ et $\mu = \sum_{i=1}^k \lambda_i \neq 0$ on note \text{$H = bar((A_1, \lambda_1), \ldots, (A_k, \lambda_k))$} alors
        $$
        G = bar((H, \mu), (A_{k+1}, \lambda_{k+1}), \ldots, (A_n, \lambda_n))
        $$
        c'est l'associativité du barycentre.
    \end{enumerate}
\end{proposition}

\begin{demonstration}
    \begin{enumerate}[label=\roman*.]
        \item.
        \item.
        \item Soit $\mu_i = \alpha \lambda_i$ pour $i \in \llbracket 1, n \rrbracket$
        On a
        \begin{align*}
            bar((A_1, \mu_1), \ldots, (A_n, \mu_n)) &\iff \sum_{i=1}^n \mu_i \overrightarrow{G A_i} = \vec{0} \\
            &\iff \sum_{i=1}^n \alpha \lambda_i \overrightarrow{G A_i} = \vec{0} \\
            &\iff \alpha \sum_{i=1}^n \lambda_i \overrightarrow{G A_i} = \vec{0} \\
            &\iff \sum_{i=1}^n \lambda_i \overrightarrow{G A_i} = \vec{0} \\
        \end{align*}
        Donc le barycentre ne change pas. Le poids total est bien $\alpha \lambda$.
        \item 
        On note $G' = bar((H, \mu), (A_{k+1}, \lambda_{k+1}), \ldots, (A_n, \lambda_n))$
        ainsi on a
        \begin{align*}
            \mu \overrightarrow{G'H} + \sum_{i=k+1}^n \lambda_i \overrightarrow{G'A_i} = \vec{0} &\iff \sum_{i=1}^k \lambda_i \overrightarrow{G'H} + \sum_{i=k+1}^n \lambda_i \overrightarrow{G'A_i} = \vec{0} \\
            &\iff \sum_{i=1}^k \lambda_i \overrightarrow{G' A_i} + \sum_{i=1}^k \lambda_i \overrightarrow{A_i H} + \sum_{i=k+1}^n \lambda_i \overrightarrow{G'A_i} = \vec{0} \\
            &\iff \sum_{i=1}^n \lambda_i \overrightarrow{G'A_i} + \sum_{i=1}^k \lambda_i \overrightarrow{A_i H} = \vec{0} \\
            &\iff \sum_{i=1}^n \lambda_i \overrightarrow{G'A_i} = \vec{0} \\
        \end{align*}
        Or $G$ est l'unique point qui vérifie cette équation, donc $G' = G$.
    \end{enumerate}
\end{demonstration}

\begin{definition}
    Si tous les poids sont égaux le barycentre s'appelle l'isobarycentre.
\end{definition}

\clearpage

\begin{theoreme}{}{}
    \begin{enumerate}
        \item Soit $\F \subset \E$ tel que $\F \neq \emptyset$. Alors $\F$ est un sous-espace affine de $\E$ si et seulement si tout barycentre d'un système de points pondérés de
        $\F$ est encore un point de $\F$.

        \item Soit $X \in \E$ tel que $X \neq \emptyset$ alors $Aff(X) = A + Vect(\overrightarrow{AM} | M \in X)$ pour un $A \in X$ est l'ensemble
        des barycentres $bar(\mathcal{A})$ où $\mathcal{A}$ parcourt l'ensemble des points pondérés de $X$.
    \end{enumerate}
\end{theoreme}

\begin{demonstration}
    \begin{enumerate}
        \item 
        \begin{description}
            \item[Sens direct.] 
            Si $\F$ est un sous-espace affine de $\E$ et si $G = bar(\mathcal{A})$ où $\mathcal{A} = ((A_1, \alpha_1), \ldots, (A_n, \alpha_n))$ et $A_i \in \F$,
            Alors $\overrightarrow{A_1 G} = \dfrac{1}{\alpha} \sum_{i=1}^n \alpha_i \overrightarrow{A_1 A_i} = \dfrac{1}{\alpha} \sum_{i=2}^n \alpha_i \overrightarrow{A_1 A_i}$
            d'où $G = A_1 + \sum_{i=2}^n \dfrac{\alpha_i}{\alpha} \overrightarrow{A_1 A_i}$ mais $\overrightarrow{A_1 A_i} \in \overrightarrow{\F}$, $\dfrac{a_i}{\alpha} \in \R$
            et $\sum_{i=2}^n \dfrac{\alpha_i}{\alpha} \in \overrightarrow{\F}$ donc $G \in \F$ car il s'écrit comme $A_1 + \vec{v}$ avec $A_1 \in \F$ et $\vec{v} \in \overrightarrow{\F}$.
            \item[Sens réciproque.]
            On suppose que $\F$ est stable par la prise de barycentre et on veut montrer que $\F$ est un sous-espace affine.
            Comme $\F \neq \emptyset$, il existe $A \in \F$.

            On considère $\overrightarrow{V} = \{ \overrightarrow{AM} | M \in \F \}$. Montrons que $\overrightarrow{V}$ est un sous-espace vectoriel de $\overrightarrow{E}$.
            On peut écrire $\overrightarrow{AA} = \overrightarrow{0} \in \overrightarrow{V}$ donc $\overrightarrow{V} \neq \emptyset$.
            Soient $\overrightarrow{AM}, \overrightarrow{AN} \in \overrightarrow{V}$ et $\lambda\in \R$.
            On a 
            \begin{align*}
                \overrightarrow{AM} + \lambda \overrightarrow{AN} = \overrightarrow{AP} &\iff \overrightarrow{AP} = \overrightarrow{AP} + \overrightarrow{PM} + \lambda (\overrightarrow{AP} + \overrightarrow{PN}) \\
                &\iff \vec{0} = \overrightarrow{PM} + \lambda \overrightarrow{PN} - \lambda \overrightarrow{PA} \\
                &\iff P = bar((M, 1), (N, \lambda), (A, -\lambda))
            \end{align*}
            où $P = bar((M, \lambda), (N, 1 - \lambda), (A, 0))$.
            Par hypothèse, $P \in \F$ donc $\overrightarrow{AP} \in \overrightarrow{V}$.
            Ainsi, $\overrightarrow{V}$ est un sous-espace vectoriel de $\overrightarrow{E}$.
            $\F = A + Vect(\overrightarrow{AM} | M \in \F)$ est bien un sous-espace affine avec $\overrightarrow{\F} = \overrightarrow{V}$.
        \end{description}
    \end{enumerate}
\end{demonstration}

\subsection{Repère barycentrique}

\begin{definition}[(Repère affine)]
    Soit $\E$ un espace affine de dimension $n$ et de direction $\vec{\E}$. Un \textbf{repère affine} de $\E$ est une famille $(A_0, A_1, \ldots, A_n)$ de points de $\E$ tels que
    $(\overrightarrow{A_0 A_1}, \ldots, \overrightarrow{A_0 A_n})$ soit une base de $\vec{\E}$.

    Ainsi pour tout $M \in \E$, il existe des réels uniques $(\alpha_1, \ldots, \alpha_n)$ tels que
    $$
    \overrightarrow{A_0 M} = \sum_{i=1}^n \alpha_i \overrightarrow{A_0 A_i}
    $$
    Les réels $\alpha_i$ sont appelés les \textbf{coordonnées barycentriques} de $M$ dans le repère affine $(A_0, A_1, \ldots, A_n)$.
\end{definition}

\begin{methode}
    \textbf{Passage entre coordonnées barycentriques et coordonnées dans un repère affine}
    \begin{align*}
        \sum_{i=0}^n \alpha_i \overrightarrow{M A_i} = \vec{0} &\iff \sum_{i=0}^n \alpha_i (\overrightarrow{MA_0} + \overrightarrow{A_0 A_i}) = \vec{0} \\
        &\iff \alpha_i \overrightarrow{M A_0} + \sum_{i=0}^n \alpha_i \overrightarrow{A_0 A_i} = \vec{0} \\
        &\iff \sum_{i=0}^n \alpha_i \overrightarrow{A_0 A_i} = (\sum_{i=0}^n \alpha_i) \overrightarrow{A_0 M} \\
        &\iff \overrightarrow{A_0 M} = \sum_{i=1}^n \alpha_i \overrightarrow{A_0 A_i}
    \end{align*}
\end{methode}

\clearpage

\begin{proposition}{}{}
    Soient $\E$ et $\mathcal{F}$ deux espaces affines. Une application $f: \E \to \mathcal{F}$ est affine si et seulement si elle préserve les barycentres.
    
    \vspace{0.5cm}
    Autrement dit, pour tout système de points pondérés \text{$\mathcal{A} = ((A_1, \lambda_1), \dots, (A_n, \lambda_n))$} avec $\sum_{i=1}^n \lambda_i \neq 0$, on a
    $$
    f(bar(\mathcal{A})) = bar(f(\mathcal{A}))
    $$
    où $f(\mathcal{A}) := ((f(A_1), \lambda_1), \dots, (f(A_n), \lambda_n))$.
\end{proposition}

\begin{demonstration}
    \begin{description}
        \item[Sens direct.] %%%%%%%%%%%%%%%%%%%%%%%%
        Supposons $f$ affine, en notant $G = bar(\mathcal{A})$, pour tout $A \in \E$, on a
        \begin{align*}
            f(G) &= f(A + \overrightarrow{AG}) \\
        &= f(A + \sum_{i=1}^n \lambda_i \overrightarrow{A A_i}) \\
        &= f(A) + \vec{f}(\sum_{i=1}^n \lambda_i \overrightarrow{A A_i}) \\
        &= f(A) + \sum_{i=1}^n \lambda_i \vec{f}(\overrightarrow{A A_i}) \\
        &= f(A) + \sum_{i=1}^n \lambda_i \overrightarrow{f(A) f(A_i)} \\
        \end{align*}
        D'où
        $$
        \overrightarrow{f(A) f(G)} = \sum_{i=1}^n \lambda_i \overrightarrow{f(A) f(A_i)}
        $$
        Ainsi $f(G)$ est le barycentre de $((f(A_1), \lambda_1), \dots, (f(A_n), \lambda_n))$.
        
        \item[Sens réciproque.] %%%%%%%%%%%%%%%%%%%%%%
        Supposons que $f$ préserve les barycentres. Montrons que $f$ est affine.

        Soit $A \in \E$ et $A' \in \E$ tels que $A' = f(A)$.

        On fixe un point $O \in \E$ et on définit une application $\funcdef{\vec{f}}{\vec{\E}}{\vec{\mathcal{F}}}{\vec{u} = \overrightarrow{OM}}{\overrightarrow{f(O)f(M)}}$

        Soit $\vec{u} \in \vec{\E}$, il existe un unique $M \in \E$ tel que $\vec{u} = \overrightarrow{OM}$.

        \begin{description}
            \item[Homogénéité.] %%%%%%%%%%%%%%%%%%%%%%%
            Soit $P \in \E$ et $\lambda \in \R$ tels que $\overrightarrow{OP} = \lambda \vec{u} = \lambda \overrightarrow{OM}$.
            
            On peut voir $P$ comme un barycentre de $(O, \alpha), (M, \beta)$, on peut trouver $\alpha$ et $\beta$ avec
            \begin{align*}
                \overrightarrow{OP} &= \dfrac{1}{\alpha + \beta} (\alpha \overrightarrow{OO} + \beta \overrightarrow{OM}) \\
                &= \dfrac{\beta}{\alpha + \beta} \overrightarrow{OM} \implies \lambda = \dfrac{\beta}{\alpha + \beta} \\
                &\iff \beta = \lambda (\alpha + \beta) \iff \beta - \lambda \beta = \lambda \alpha \iff \beta (1 - \lambda) = \lambda \alpha \\
                &\iff \dfrac{\beta}{\alpha} = \dfrac{\lambda}{1 - \lambda} \iff \begin{cases}
                    \alpha = 1 - \lambda \\
                    \beta = \lambda
                \end{cases}
            \end{align*}

            Ainsi,
            \begin{align*}
                \vec{f}(\lambda \vec{u}) &= \vec{f}(\overrightarrow{OP}) \\
                &= \overrightarrow{f(O) f(P)} \\
                &= \overrightarrow{f(O) f(bar((O, 1 - \lambda), (M, \lambda)))} \\
                &= \overrightarrow{f(O) bar((f(O), 1 - \lambda), (f(M), \lambda))} \quad \text{par hypothèse} \\
                &= \dfrac{1}{(1 - \lambda) + \lambda} ((1 - \lambda) \overrightarrow{f(O) f(O)} + \lambda \overrightarrow{f(O) f(M)}) \\
                &= \lambda \overrightarrow{f(O) f(M)} \\
                &= \lambda \vec{f}(\vec{u})
            \end{align*}

            \item[Linéarité.] %%%%%%%%%%%%%%%%%%%%%%%
            Soit $\vec{u}, \vec{v} \in \vec{\E}$, il existe $M, N \in \E$ tels que $\vec{u} = \overrightarrow{OM}$ et $\vec{v} = \overrightarrow{ON}$

            Soit $I \in \E$ le barycentre de $(M, 1), (N, 1)$ (le milieu quoi!).

            ainsi on a,
            $$
            \overrightarrow{OI} = \dfrac{\overrightarrow{OM} + \overrightarrow{ON}}{2} = \dfrac{\vec{u} + \vec{v}}{2} \implies \vec{f}(\overrightarrow{OI}) = \vec{f}\left(\dfrac{\vec{u} + \vec{v}}{2}\right)
            $$
            ainsi,
            \begin{align*}
                \vec{f}(\vec{u} + \vec{v}) &= \vec{f}(2 \overrightarrow{OI}) \\
                &= 2 \vec{f}(\overrightarrow{OI}) \\
                &= 2 \vec{f}\left(\dfrac{\vec{u} + \vec{v}}{2}\right) \\
                &= \vec{f}(\vec{u}) + \vec{f}(\vec{v})
            \end{align*}
        \end{description}
        Donc $\vec{f}$ est linéaire et $f$ est affine.
    \end{description}
\end{demonstration}

\begin{proposition}{}{}
    Soient $\E$ et $\mathcal{F}$ deux espaces affines de dimensions finies $n$.

    Soit $(A_0, \dots, A_n)$ un repère affine/barycentrique de $\E$ et \text{$(B_0, \dots, B_n) \in \mathcal{F}^n$}.
    
    \vspace{0.5cm}
    Alors, il existe une unique application affine $f: \E \to \mathcal{F}$ telle que $f(A_i) = B_i$ pour tout $i \in \llbracket 0, n \rrbracket$ et $f$ est un isomorphisme si et seulement si
    $(B_0, \dots, B_n)$ est un repère affine/barycentrique de $\mathcal{F}$.
\end{proposition}

\begin{demonstration}
    \begin{description}
        \item[Unicité.] 
        Soient $f$ et $g$ deux applications affines telles que $f(A_i) = B_i$ et $g(A_i) = B_i$.
        
        Calculons, pour $M \in \E$,
        \begin{align*}
            f(M) &= f(A_0 + \overrightarrow{A_0 M}) \\
            &= f\left(A_0 + \sum_{i=1}^n \lambda_i \overrightarrow{A_0 A_i}\right) \\
            &= f(A_0) + \sum_{i=1}^n \lambda_i \vec{f}(\overrightarrow{A_0 A_i}) \\
            &= B_0 + \sum_{i=1}^n \lambda_i \overrightarrow{B_0 B_i} \\
            &= g(A_0) + \sum_{i=1}^n \lambda_i \vec{g}(\overrightarrow{A_0 A_i}) \\
            &= g\left(A_0 + \sum_{i=1}^n \lambda_i \overrightarrow{A_0 A_i}\right) \\
            &= g(M)
        \end{align*}
        Donc $f = g$ et l'unicité est montrée.

        \item[Bijectivité.]
        On définit une application affine $g: \mathcal{F} \to \E$ telle que pour tout $i \in \llbracket 0, n \rrbracket$, $g(B_i) = A_i$.
        
        Soit $M \in \E$, on a
        \begin{align*}
            g(f(M)) &= g\left(B_0 + \sum_{i=1}^n \lambda_i \overrightarrow{B_0 B_i}\right) \\
            &= g(B_0) + \sum_{i=1}^n \lambda_i \vec{g}(\overrightarrow{B_0 B_i}) \\
            &= A_0 + \sum_{i=1}^n \lambda_i \overrightarrow{A_0 A_i} \\
            &= M
        \end{align*}
        Ainsi $g \circ f = id_{\E}$.

        De même, pour $N \in \mathcal{F}$, on a
        \begin{align*}
            f(g(N)) &= f\left(A_0 + \sum_{i=1}^n \mu_i \overrightarrow{A_0 A_i}\right) \\
            &= f(A_0) + \sum_{i=1}^n \mu_i \vec{f}(\overrightarrow{A_0 A_i}) \\
            &= B_0 + \sum_{i=1}^n \mu_i \overrightarrow{B_0 B_i} \\
            &= N
        \end{align*}
        Ainsi $f \circ g = id_{\mathcal{F}}$.

        Donc $f$ est un isomorphisme.
\end{description}
\end{demonstration}

\clearpage

\section{Le triangle}

\subsection{Droites remarquables d'un triangle}

\begin{definition}
    Soit $ABC$ un triangle non plat. On appelle \textbf{cercle circonscrit} du triangle $ABC$ le cercle qui passe par les trois sommets $A$, $B$ et $C$.
\end{definition}

\begin{definition}
    Soit $ABC$ un triangle. On appelle \textbf{mediatrice} du segment $[BC]$ la droite qui est perpendiculaire à la droite $(BC)$ et qui passe par le milieu de $[BC]$.
\end{definition}

\begin{proposition}{}{}
    La mediatrice du segment $[BC]$ est l'ensemble des points équidistants de $B$ et de $C$.
\end{proposition}

\begin{theoreme}{}{}
    Les trois mediatrices d'un triangle sont concourantes. Le point de concours est appelé le \textbf{centre du cercle circonscrit} du triangle.
    
    Il est souvent noté $O$.
\end{theoreme}

\tikzexternalenable
\begin{figure}[H]
    \centering
    \resizebox{0.44\textwidth}{!}{%
    \begin{tikzpicture}
    \tikzset{
        mark 1/.style={postaction={decorate, decoration={markings, mark=at position 0.5 with {\draw[thick] (0,-3.5pt) -- (0,3.5pt);}}}},
        mark 2/.style={postaction={decorate, decoration={markings, mark=at position 0.5 with {\draw[thick] (-1.5pt,-3.5pt) -- (-1.5pt,3.5pt); \draw[thick] (1.5pt,-3.5pt) -- (1.5pt,3.5pt);}}}},
        mark 3/.style={postaction={decorate, decoration={markings, mark=at position 0.5 with {\draw[thick] (-2.5pt,-3.5pt) -- (-2.5pt,3.5pt); \draw[thick] (0,-3.5pt) -- (0,3.5pt); \draw[thick] (2.5pt,-3.5pt) -- (2.5pt,3.5pt);}}}}
    }

    \coordinate (O) at (0,0);
    \def\R{4}

    \coordinate (A) at (110:\R);
    \coordinate (B) at (210:\R);
    \coordinate (C) at (330:\R);

    \coordinate (Mc) at ($(A)!0.5!(B)$); % Milieu de AB
    \coordinate (Ma) at ($(B)!0.5!(C)$); % Milieu de BC
    \coordinate (Mb) at ($(A)!0.5!(C)$); % Milieu de AC

    \draw[red, thick] (O) circle (\R);

    \draw[blue, thin] ($(O)!-0.3!(Mc)$) -- ($(O)!1.3!(Mc)$);
    \draw[blue, thin] ($(O)!-0.3!(Ma)$) -- ($(O)!1.3!(Ma)$);
    \draw[blue, thin] ($(O)!-0.3!(Mb)$) -- ($(O)!1.3!(Mb)$);

    \draw[black, very thick] (A) -- (B) -- (C) -- cycle;

    \pic [draw, angle radius=3.5mm] {right angle = O--Mc--A};
    \pic [draw, angle radius=3.5mm] {right angle = O--Ma--C};
    \pic [draw, angle radius=3.5mm] {right angle = O--Mb--A};

    \draw[mark 1] (A) -- (Mc); \draw[mark 1] (Mc) -- (B);
    \draw[mark 2] (B) -- (Ma); \draw[mark 2] (Ma) -- (C);
    \draw[mark 3] (A) -- (Mb); \draw[mark 3] (Mb) -- (C);

    \foreach \p in {A,B,C,O}
        \filldraw[fill=white, draw=black, thick] (\p) circle (2.5pt);

    \node[above=3pt] at (A) {A};
    \node[below left=3pt] at (B) {B};
    \node[below right=3pt] at (C) {C};
    \node[below right=0pt, yshift=-2pt] at (O) {O};

\end{tikzpicture}
    }
    \caption{Cercle circonscrit (en rouge) et médiatrices (en bleu) d'un triangle $ABC$. Les trois médiatrices se coupent en un point $O$ qui est le centre du cercle circonscrit.}
\end{figure}
\tikzexternaldisable

\begin{definition}
    Soit $ABC$ un triangle non plat. On appelle \textbf{cercle inscrit} du triangle $ABC$ le cercle qui est tangent aux trois côtés du triangle.
\end{definition}

\begin{definition}
    Soit $ABC$ un triangle. On appelle \textbf{bissectrice} issue de $A$ la droite qui passe par $A$ et qui partage l'angle $\widehat{BAC}$ en deux angles de même mesure.
\end{definition}

\begin{proposition}{}{}
    La bissectrice issue de $A$ est l'ensemble des points équidistants des droites $(AB)$ et $(AC)$.
\end{proposition}

\begin{theoreme}{}{}
    Les trois bissectrices d'un triangle sont concourantes. Le point de concours est appelé le \textbf{centre du cercle inscrit} du triangle.
    
    Il est souvent noté $I$.
\end{theoreme}

\tikzexternalenable
\begin{figure}[H]
    \centering
    \resizebox{0.6\textwidth}{!}{%
    \begin{tikzpicture}

    \tikzset{
        mark segment/.style={postaction={decorate, decoration={markings, mark=at position 0.5 with {\draw[thick] (-1.5pt,-3.5pt) -- (-1.5pt,3.5pt); \draw[thick] (1.5pt,-3.5pt) -- (1.5pt,3.5pt);}}}},
        tick 1/.style={
            decoration={markings, mark=at position 0.5 with {\draw[thick] (0,-3.5pt) -- (0,3.5pt);}},
            postaction={decorate}
        },
        tick 2/.style={
            decoration={markings, mark=at position 0.5 with {\draw[thick] (-1.5pt,-3.5pt) -- (-1.5pt,3.5pt); \draw[thick] (1.5pt,-3.5pt) -- (1.5pt,3.5pt);}},
            postaction={decorate}
        },
        tick 3/.style={
            decoration={markings, mark=at position 0.5 with {\draw[thick] (-2.5pt,-3.5pt) -- (-2.5pt,3.5pt); \draw[thick] (0,-3.5pt) -- (0,3.5pt); \draw[thick] (2.5pt,-3.5pt) -- (2.5pt,3.5pt);}},
            postaction={decorate}
        }
    }

    % --- Définition des points principaux ---
    \coordinate (A) at (2.5, 6);
    \coordinate (B) at (0, 0);
    \coordinate (C) at (7, 0);
    
    % Centre du cercle inscrit (I) et rayon
    \coordinate (I) at (3, 2); 
    \def\r{2} 

    % --- Tracé du cercle inscrit ---
    \draw[red, thick] (I) circle (\r);

    % --- Calcul des intersections pour les bissectrices ---
    \path[name path=sidea] (B) -- (C);
    \path[name path=sideb] (A) -- (C);
    \path[name path=sidec] (A) -- (B);

    \path[name path=lineA] (A) -- ($(A)!2!(I)$);
    \path[name path=lineB] (B) -- ($(B)!2!(I)$);
    \path[name path=lineC] (C) -- ($(C)!2!(I)$);

    \path [name intersections={of=lineA and sidea, by=IntA}];
    \path [name intersections={of=lineB and sideb, by=IntB}];
    \path [name intersections={of=lineC and sidec, by=IntC}];

    % --- Construction de l'équidistance au sommet C ---
    % 1. On place un point P sur la bissectrice CI (à 60% de la distance)
    \coordinate (P) at ($(C)!0.4!(I)$);
    
    % 2. On crée une ligne perpendiculaire à CI passant par P
    \coordinate (Pdir1) at ($(P)!2cm!90:(I)$);
    \coordinate (Pdir2) at ($(P)!2cm!-90:(I)$);
    \path[name path=perpC] (Pdir1) -- (Pdir2);
    
    % 3. On trouve les intersections D et E avec les côtés AC et BC
    \path [name intersections={of=perpC and sideb, by=D}]; 
    \path [name intersections={of=perpC and sidea, by=E}]; 

    % --- Tracé du triangle et des bissectrices ---
    \draw[black, very thick] (A) -- (B) -- (C) -- cycle;

    \draw[blue, thin] (A) -- ($(A)!1.15!(IntA)$);
    \draw[blue, thin] (B) -- ($(B)!1.15!(IntB)$);
    \draw[blue, thin] (C) -- ($(C)!1.15!(IntC)$);

    % --- Dessin de la propriété au sommet C ---
    \draw[orange, thick] (D) -- (E); % Le segment perpendiculaire
    \pic [draw, angle radius=2.5mm] {right angle = I--P--D}; % L'angle droit
    \draw[mark segment] (D) -- (P); % Marque d'égalité PD
    \draw[mark segment] (P) -- (E); % Marque d'égalité PE

    % --- Nouvelles marques des angles (façon "arc barré") ---
    % Sommet B (1 trait)
    \pic [draw, angle radius=8mm] {angle = C--B--I};
    \pic [draw, angle radius=6mm] {angle = I--B--A};

    % Sommet C (2 traits)
    \pic [draw, angle radius=8mm] {angle = A--C--I};
    \pic [draw, angle radius=7mm] {angle = A--C--I};
    \pic [draw, angle radius=6mm] {angle = I--C--B};
    \pic [draw, angle radius=5mm] {angle = I--C--B};

    % Sommet A (3 traits)
    \pic [draw, angle radius=9mm] {angle = B--A--I};
    \pic [draw, angle radius=8mm] {angle = B--A--I};
    \pic [draw, angle radius=7mm] {angle = B--A--I};
    \pic [draw, angle radius=6mm] {angle = I--A--C};
    \pic [draw, angle radius=5mm] {angle = I--A--C};
    \pic [draw, angle radius=4mm] {angle = I--A--C};

    % --- Ajout des points et labels ---
    \foreach \p in {A,B,C,I}
        \filldraw[fill=white, draw=black, thick] (\p) circle (2pt);

    \node[above=3pt] at (A) {A};
    \node[below left=3pt] at (B) {B};
    \node[below right=3pt] at (C) {C};
    \node[above left=1pt, yshift=2pt] at (I) {I};


\end{tikzpicture}
    }
    \caption{Cercle inscrit (en rouge) et bissectrices (en bleu) d'un triangle $ABC$. Les trois bissectrices se coupent en un point $I$ qui est le centre du cercle inscrit.}
\end{figure}
\tikzexternaldisable

\begin{definition}
    Soit $ABC$ un triangle. On appelle \textbf{médiane} issue de $A$ la droite qui passe par $A$ et par le milieu de $[BC]$.
\end{definition}

\begin{proposition}{}{}
    Le centre de gravité d'un triangle est situé à $\frac{2}{3}$ de la médiane issue de chaque sommet à partir du sommet.
\end{proposition}

\begin{theoreme}{}{}
    Les trois médianes d'un triangle sont concourantes. Le point de concours est appelé le \textbf{centre de gravité} du triangle.
    
    Il est souvent noté $G$.
\end{theoreme}

\tikzexternalenable
\begin{figure}[H]
    \centering
    \resizebox{0.6\textwidth}{!}{%
    \begin{tikzpicture}

    % --- Styles robustes pour les marques de segments ---
    \tikzset{
        % Permet de dessiner 1, 2 ou 3 petits traits au milieu du segment
        mark 1/.style={postaction={decorate, decoration={markings, mark=at position 0.5 with {\draw[thick] (0,-3.5pt) -- (0,3.5pt);}}}},
        mark 2/.style={postaction={decorate, decoration={markings, mark=at position 0.5 with {\draw[thick] (-1.5pt,-3.5pt) -- (-1.5pt,3.5pt); \draw[thick] (1.5pt,-3.5pt) -- (1.5pt,3.5pt);}}}},
        mark 3/.style={postaction={decorate, decoration={markings, mark=at position 0.5 with {\draw[thick] (-2.5pt,-3.5pt) -- (-2.5pt,3.5pt); \draw[thick] (0,-3.5pt) -- (0,3.5pt); \draw[thick] (2.5pt,-3.5pt) -- (2.5pt,3.5pt);}}}}
    }

    % --- Définition des sommets du triangle ---
    \coordinate (A) at (1.5, 5);
    \coordinate (B) at (0, 0);
    \coordinate (C) at (7, 0);

    % --- Calcul des milieux des côtés ---
    \coordinate (Aprime) at ($(B)!0.5!(C)$); % Milieu de BC (opposé à A)
    \coordinate (Bprime) at ($(A)!0.5!(C)$); % Milieu de AC (opposé à B)
    \coordinate (Cprime) at ($(A)!0.5!(B)$); % Milieu de AB (opposé à C)

    % --- Calcul du Centre de Gravité (G) ---
    % TikZ sait calculer les barycentres nativement, c'est parfait pour G !
    \coordinate (G) at (barycentric cs:A=1,B=1,C=1);

    % --- Tracé du triangle ---
    \draw[black, very thick] (A) -- (B) -- (C) -- cycle;

    % --- Tracé des médianes ---
    \draw[blue, thin] (A) -- (Aprime);
    \draw[blue, thin] (B) -- (Bprime);
    \draw[blue, thin] (C) -- (Cprime);

    % --- Ajout des marques de segments (équidistance) ---
    % Côté BC (coupé par A') -> 1 trait
    \draw[mark 1] (B) -- (Aprime); 
    \draw[mark 1] (Aprime) -- (C);
    
    % Côté AC (coupé par B') -> 2 traits
    \draw[mark 2] (A) -- (Bprime); 
    \draw[mark 2] (Bprime) -- (C);

    % Côté AB (coupé par C') -> 3 traits
    \draw[mark 3] (A) -- (Cprime); 
    \draw[mark 3] (Cprime) -- (B);

    % --- Ajout des points et labels ---
    \foreach \p in {A,B,C,G,Aprime,Bprime,Cprime}
        \filldraw[fill=white, draw=black, thick] (\p) circle (2pt);

    \node[above=3pt] at (A) {A};
    \node[below left=3pt] at (B) {B};
    \node[below right=3pt] at (C) {C};
    
    % Labels des milieux
    \node[below=3pt] at (Aprime) {A'};
    \node[above right=2pt] at (Bprime) {B'};
    \node[above left=2pt] at (Cprime) {C'};

    % Label du centre de gravité
    \node[above right=0pt, yshift=2pt] at (G) {G};

\end{tikzpicture}
    }
    \caption{Médianes (en bleu) d'un triangle $ABC$. Les trois médianes se coupent en un point $G$ qui est le centre de gravité du triangle.}
\end{figure}
\tikzexternaldisable

\begin{definition}
    Soit $ABC$ un triangle. On appelle \textbf{hauteur} issue de $A$ la droite qui passe par $A$ et qui est perpendiculaire à la droite $(BC)$.
\end{definition}

\begin{theoreme}{}{}
    Les trois hauteurs d'un triangle sont concourantes. Le point de concours est appelé l'\textbf{orthocentre} du triangle.
    
    Il est souvent noté $H$.
\end{theoreme}

\tikzexternalenable
\begin{figure}[H]
    \centering
    \resizebox{0.6\textwidth}{!}{%
    \begin{tikzpicture}

    % --- Définition des sommets du triangle ---
    % On choisit un triangle acutangle (tous les angles < 90°) 
    % pour que l'orthocentre soit bien à l'intérieur de la figure.
    \coordinate (A) at (2, 5);
    \coordinate (B) at (0, 0);
    \coordinate (C) at (7, 0);

    % --- Calcul des pieds des hauteurs par projection orthogonale ---
    % La syntaxe magique de TikZ : $(B)!(A)!(C)$ projette le point A sur la droite (BC)
    \coordinate (Ha) at ($(B)!(A)!(C)$); 
    \coordinate (Hb) at ($(A)!(B)!(C)$); 
    \coordinate (Hc) at ($(A)!(C)!(B)$);

    % --- Calcul de l'Orthocentre (H) ---
    % On le trouve en cherchant l'intersection de deux hauteurs
    \path[name path=altA] (A) -- (Ha);
    \path[name path=altB] (B) -- (Hb);
    \path [name intersections={of=altA and altB, by=H}];

    % --- Tracé du triangle ---
    \draw[black, very thick] (A) -- (B) -- (C) -- cycle;

    % --- Tracé des hauteurs ---
    \draw[blue, thin] (A) -- (Ha);
    \draw[blue, thin] (B) -- (Hb);
    \draw[blue, thin] (C) -- (Hc);

    % --- Ajout des symboles d'angles droits ---
    \pic [draw, angle radius=3mm] {right angle = A--Ha--C};
    \pic [draw, angle radius=3mm] {right angle = B--Hb--C};
    \pic [draw, angle radius=3mm] {right angle = C--Hc--B};

    % --- Ajout des points et labels ---
    \foreach \p in {A,B,C,H}
        \filldraw[fill=white, draw=black, thick] (\p) circle (2pt);

    \node[above=3pt] at (A) {A};
    \node[below left=3pt] at (B) {B};
    \node[below right=3pt] at (C) {C};
    
    % Labels des pieds des hauteurs
    \node[below=3pt] at (Ha) {$H_A$};
    \node[above right=1pt] at (Hb) {$H_B$};
    \node[above left=1pt] at (Hc) {$H_C$};

    % Label de l'orthocentre
    \node[above left=1pt, yshift=2pt] at (H) {H};

\end{tikzpicture}
    }
    \caption{Hauteurs (en bleu) d'un triangle $ABC$. Les trois hauteurs se coupent en un point $H$ qui est l'orthocentre du triangle.}
\end{figure}
\tikzexternaldisable

\subsection{Les propriétés élémentaires}

\clearpage

\section{Les angles}

\subsection{My name is BOND, James BOND}

\begin{definition}[(Orientation d'un espace vectoriel)]
    Soit $E$ un $\R$-espace vectoriel de dimension finie.

    \begin{itemize}
        \item Deux bases $\mathcal{B}$ et $\mathcal{B}'$ de $E$ ont la \textbf{même orientation} si la matrice de passage de $\mathcal{B}$ à $\mathcal{B}'$ a un déterminant strictement positif.
        \item La relation "avoir la même orientation" est une relation d'équivalence sur l'ensemble des bases de $E$. Il y a exactement deux classes d'équivalence, appelées les \textbf{orientations} de $E$.
        \item Orienter $E$ consiste à choisir arbitrairement une orientation de $E$, c'est à dire une classe d'équivalence de bases de $E$.
        Les bases appartenant à l'orientation choisie sont dites \textbf{bases directes}, les autres sont dites \textbf{bases indirectes}.
    \end{itemize}
\end{definition}

\begin{definition}[(Base orthonormée directe ou BOND)]
    Soit $E$ un espace euclidien orienté de dimension finie.

    Une base $\mathcal{B}$ de $E$ est une \textbf{base orthonormée directe} (BOND) si elle verifie les conditions
    \begin{itemize}
        \item $\mathcal{B}$ est une base orthonormée de $E$.
        \item $\mathcal{B}$ est une base directe de $E$.
    \end{itemize}
\end{definition}

\subsection{Angle de vecteurs}

\begin{theoreme}{}{angle_vecteurs}
    Soit $\mathcal{P}$ un plan vectoriel euclidien orienté.

    Soient $u, v \in \mathcal{P}$ tels que $\norm{u} = \norm{v} = 1$.

    \vspace{1cm}
    En notant $(a, b)$ les coordonnées de $v$ dans la BOND de premier vecteur $u$,
    il existe un unique $\theta \in \R$ modulo $2\pi$ tel que $a = \cos(\theta)$ et $b = \sin(\theta)$.
\end{theoreme}

\begin{demonstration}
    Il existe un unique $u^\perp \in \mathcal{P}$ tel que $\mathcal{B} = (u, u^\perp)$ soit une BOND de $\mathcal{P}$.

    On considère l'unique rotation vectorielle $r$ telle que $r(u) = v$.

    Ainsi on peut écrire
    $$
    Mat_{\mathcal{B}}(r) = \begin{pmatrix}
        a & -b \\
        b & a
    \end{pmatrix}
    $$

    Une rotation vectorielle conserve l'orientation et les distances, donc $det(r) = 1$ et
    $$
    a^2 + b^2 = 1
    $$
    Il existe un unique $\theta \in \R$ modulo $2\pi$ tel que $a = \cos(\theta)$ et $b = \sin(\theta)$.
\end{demonstration}

\begin{definition}[(Angle de vecteurs)]
    Soit $\mathcal{P}$ un plan vectoriel euclidien orienté.

    Soient $u, v \in \mathcal{P}\setminus\{0\}$.

    On appelle l'\textbf{angle orienté} de $u$ à $v$, noté $(\widehat{\overrightarrow{u}, \overrightarrow{v}})$,
    le réel $\theta$ défini par le théorème~\ref{theoreme:angle_vecteurs} appliqué à $\frac{u}{\norm{u}}$ et $\frac{v}{\norm{v}}$.
\end{definition}

\begin{remarque}
    L'ensemble des angles orientés de vecteurs du plan est isomorphe à $\R / 2\pi\Z$.
    C'est un groupe abélien pour l'addition.
\end{remarque}

\subsection{Angle de droites}

\begin{proposition}{}{}
    Soient $D$ et $D'$ deux droites d'un plan vectoriel euclidien orienté.

    Si $u_1$ et $u_2$ sont deux vecteurs directeurs de $D$, et $v_1$ et $v_2$ sont deux vecteurs directeurs de $D'$, alors
    $$
    (\widehat{\overrightarrow{u_1}, \overrightarrow{v_1}}) \equiv (\widehat{\overrightarrow{u_2}, \overrightarrow{v_2}}) \mod \pi
    $$
\end{proposition}

\begin{definition}
    Soient $D$ et $D'$ deux droites d'un plan vectoriel euclidien orienté et $u$ et $v$ des vecteurs directeurs de $D$ et $D'$ respectivement.

    On définit l'\textbf{angle orienté} de $D$ à $D'$ par
    $$
    (\widehat{D, D'}) = (\widehat{\overrightarrow{u}, \overrightarrow{v}}) \mod \pi
    $$
\end{definition}

\begin{remarque}
    Par cette définition, on a le resultat peu intuitif de l'égalité $\alpha = \beta$ dans la figure ci-dessous (car $\beta = \alpha + \pi$)
\end{remarque}

\tikzexternalenable
\begin{figure}[H]
    \centering
    \resizebox{0.6\textwidth}{!}{%
    \begin{tikzpicture}[>=stealth, scale=2]

% --- Droites ---
\coordinate (O) at (0,0);
\coordinate (A) at (1.5,0);
\coordinate (B) at (1.5,1.5);
\coordinate (A_neg) at (-1.5,0);
\coordinate (B_neg) at (-1.5,-1.5);

% Droite D
\draw[thick, red] (A_neg) -- (A) node[anchor=east, pos=0.9, fill=white] {$D$};
% Droite D'
\draw[thick, blue] (B_neg) -- (B) node[anchor=south west, pos=0.7, fill=white] {$D'$};

% --- Vecteurs directeurs ---
% Vecteur u sur la demi-droite positive de D
%\draw[->, ultra thick, red!60!black] (O) -- (1,0) node[pos=0.8, yshift=6pt, scale=1.2] {$\mathbf{u}$};
% Vecteur -u sur la demi-droite négative de D
%\draw[->, ultra thick, red!60!black] (O) -- (-1,0) node[pos=0.8, yshift=6pt, scale=1.2] {$-\mathbf{u}$};
% Vecteur v sur la demi-droite positive de D'
%\draw[->, ultra thick, blue!60!black] (O) -- (1,1) node[pos=0.8, xshift=6pt, yshift=3pt, scale=1.2] {$\mathbf{v}$};

% --- Angles ---
% Angle alpha : de u vers v
\pic [draw, ->, angle radius=1.2cm, "$\alpha$", angle eccentricity=1.2, red!80!black] {angle = A--O--B};

% Angle beta : de -u vers v
% Nous traçons l'arc géométrique de la droite D' vers la droite D
\pic [draw, <-, angle radius=1.3cm, "$\beta$", angle eccentricity=1.35, blue!80!black] {angle = B--O--A_neg};

\end{tikzpicture}
    }
\end{figure}
\tikzexternaldisable

\begin{remarque}
    L'ensemble des angles orientés de droites d'un plan est isomorphe à $\R / \pi\Z$.
    C'est un groupe abélien pour l'addition.
\end{remarque}

\begin{remarque}
    On peut caractériser les positions relatives des droites
    \begin{itemize}
        \item $D$ et $D'$ sont parallèles si et seulement si $(\widehat{D, D'}) = 0$.
        \item $D$ et $D'$ sont perpendiculaires si et seulement si $(\widehat{D, D'}) = \frac{\pi}{2}$.
    \end{itemize}
\end{remarque}

\subsection{Chasles et mesures principales}

\begin{theoreme}{Relation de Chasles}{}
    Soit $\mathcal{P}$ un plan vectoriel euclidien orienté.
    \begin{enumerate}
        \item Pour tous $u, v, w \in \mathcal{P}\setminus\{0\}$, on a
        $$
        (\widehat{\overrightarrow{u}, \overrightarrow{v}}) + (\widehat{\overrightarrow{v}, \overrightarrow{w}}) \equiv (\widehat{\overrightarrow{u}, \overrightarrow{w}}) \mod 2\pi
        $$
        \item Pour tous $D, D', D''$ droites d'un plan vectoriel euclidien orienté, on a
        $$
        (\widehat{D, D'}) + (\widehat{D', D''}) \equiv (\widehat{D, D''}) \mod \pi
        $$
    \end{enumerate}
\end{theoreme}

\begin{demonstration}
    Quitte à normaliser les vecteurs, on peut supposer que $u$, $v$ et $w$ sont unitaire.

    Soit $r_1$ la rotation vectorielle telle que $r_1(u) = v$, d'angle $\theta_1 = (\widehat{\overrightarrow{u}, \overrightarrow{v}})$.
    
    Soit $r_2$ la rotation vectorielle telle que $r_2(v) = w$, d'angle $\theta_2 = (\widehat{\overrightarrow{v}, \overrightarrow{w}})$.

    En représentant les rotations vectorielles par des matrices dans la BOND de premier vecteur $u$, on trouve
    $$
    Mat_{\mathcal{B}}(r_1) = \begin{pmatrix}
        \cos(\theta_1) & -\sin(\theta_1) \\
        \sin(\theta_1) & \cos(\theta_1)
    \end{pmatrix}
    $$
    et
    $$
    Mat_{\mathcal{B}}(r_2) = \begin{pmatrix}
        \cos(\theta_2) & -\sin(\theta_2) \\
        \sin(\theta_2) & \cos(\theta_2)
    \end{pmatrix}
    $$
    Ainsi
    $$
    Mat_{\mathcal{B}}(r_2 \circ r_1) = \begin{pmatrix}
        \cos(\theta_1 + \theta_2) & -\sin(\theta_1 + \theta_2) \\
        \sin(\theta_1 + \theta_2) & \cos(\theta_1 + \theta_2)
    \end{pmatrix}
    $$
    Or $r_2 \circ r_1(u) = w$, donc $r_2 \circ r_1$ est la rotation vectorielle d'angle $(\widehat{\overrightarrow{u}, \overrightarrow{w}})$, d'où le résultat sur les vecteurs.

    Le résultat sur les droites s'en déduit en prenant des vecteurs directeurs de ces droites.
\end{demonstration}

\begin{corollaire}{}{}
    Pour tous $u, v \in \mathcal{P}\setminus\{0\}$, on a
    $$
    (\widehat{\overrightarrow{u}, \overrightarrow{v}}) \equiv -(\widehat{\overrightarrow{v}, \overrightarrow{u}}) \mod 2\pi
    $$
\end{corollaire}

\begin{definition}[(Mesure principale d'un angle)]
    Soit $\mathcal{P}$ un plan vectoriel euclidien orienté.

    \begin{itemize}
        \item Pour tous $u, v \in \mathcal{P}\setminus\{0\}$, on appelle \textbf{mesure principale} de l'angle orienté de $u$ à $v$ l'unique réel $\theta$ dans l'intervalle $]-\pi, \pi]$ tel que $$\theta \equiv (\widehat{\overrightarrow{u}, \overrightarrow{v}}) \mod 2\pi$$
        \item Pour tous $D, D'$ droites d'un plan vectoriel euclidien orienté, on appelle \textbf{mesure principale} de l'angle orienté de $D$ à $D'$ l'unique réel $\theta$ dans l'intervalle $]-\frac{\pi}{2}, \frac{\pi}{2}]$ tel que $$\theta \equiv (\widehat{D, D'}) \mod \pi$$
    \end{itemize}
\end{definition}

\subsection{Expression analytique}

\begin{theoreme}{}{}
    Soit $\mathcal{P}$ un plan vectoriel euclidien orienté.

    Soient $u, v \in \mathcal{P}\setminus\{0\}$.

    Alors
    $$
    \langle\overrightarrow{u}, \overrightarrow{v}\rangle = \norm{u}\norm{v}\cos((\widehat{\overrightarrow{u}, \overrightarrow{v}}))
    $$
    et
    $$
    det(\overrightarrow{u}, \overrightarrow{v}) = \norm{u}\norm{v}\sin((\widehat{\overrightarrow{u}, \overrightarrow{v}}))
    $$
\end{theoreme}

\begin{demonstration}
    Posons $i = \dfrac{u}{\norm{u}}$. Il existe un unique $j \in \mathcal{P}$ tel que $\mathcal{B} = (i, j)$ soit une BOND de $\mathcal{P}$.

    Les coordonnées de $u$ dans $\mathcal{B}$ sont $(\norm{u}, 0)$, et $\dfrac{v}{\norm{v}}$ a pour coordonnées $(\cos((\widehat{\overrightarrow{u}, \overrightarrow{v}})), \sin((\widehat{\overrightarrow{u}, \overrightarrow{v}})))$.

    Ainsi $v$ a pour coordonnées $(\norm{v}\cos((\widehat{\overrightarrow{u}, \overrightarrow{v}})), \norm{v}\sin((\widehat{\overrightarrow{u}, \overrightarrow{v}})))$ dans la base $\mathcal{B}$.
    Par conséquent
    $$
    \langle\overrightarrow{u}, \overrightarrow{v}\rangle = \norm{u}\norm{v}\cos((\widehat{\overrightarrow{u}, \overrightarrow{v}}))
    $$
    et
    $$
    det(\overrightarrow{u}, \overrightarrow{v}) = \begin{vmatrix}
        \norm{u} & \norm{v}\cos((\widehat{\overrightarrow{u}, \overrightarrow{v}})) \\
        0 & \norm{v}\sin((\widehat{\overrightarrow{u}, \overrightarrow{v}}))
    \end{vmatrix} = \norm{u}\norm{v}\sin((\widehat{\overrightarrow{u}, \overrightarrow{v}}))
    $$
\end{demonstration}

\subsection{Angle inscrit et cocyclicité}

\begin{lemme}{Angle du triangle isocèle}{}
    Soit $\Omega M A$ un triangle isocèle en $\Omega$, alors
    $$
    2(\widehat{\overrightarrow{MA}, \overrightarrow{M\Omega}}) \equiv (\widehat{\overrightarrow{\Omega A}, \overrightarrow{\Omega M}}) + \pi \mod 2\pi
    $$
\end{lemme}

\begin{demonstration}
    Posons $\theta_1 = (\widehat{\overrightarrow{MA}, \overrightarrow{M\Omega}})$
    et $\theta_2 = (\widehat{\overrightarrow{A \Omega}, \overrightarrow{A M}})$.

    On a
    \begin{align*}
        \langle\overrightarrow{MA}, \overrightarrow{M\Omega}\rangle &= \langle\overrightarrow{M\Omega} + \overrightarrow{\Omega A}, \overrightarrow{M\Omega}\rangle \\
        &= \norm{\overrightarrow{M\Omega}}^2 + \langle\overrightarrow{\Omega A}, \overrightarrow{M\Omega}\rangle
    \end{align*}
    et
    \begin{align*}
        \langle\overrightarrow{A \Omega}, \overrightarrow{A M}\rangle &= \langle-\overrightarrow{\Omega A}, -\overrightarrow{M A}\rangle \\
        &= (-1)(-1)\langle\overrightarrow{\Omega A}, \overrightarrow{M A}\rangle \\
        &= \langle\overrightarrow{\Omega A}, \overrightarrow{M\Omega} + \overrightarrow{\Omega A}\rangle \\
        &= \langle\overrightarrow{\Omega A}, \overrightarrow{M\Omega}\rangle + \norm{\overrightarrow{\Omega A}}^2
    \end{align*}

    Or $\norm{\overrightarrow{M\Omega}} = \norm{\overrightarrow{\Omega A}}$ car le triangle est isocèle en $\Omega$, donc $\cos(\theta_1) = \cos(\theta_2)$

    De même pour les déterminants, on trouve $\sin(\theta_1) = \sin(\theta_2)$.

    Ainsi $(\widehat{\overrightarrow{MA}, \overrightarrow{M\Omega}}) \equiv (\widehat{\overrightarrow{A \Omega}, \overrightarrow{A M}}) \mod 2\pi$.

    La somme des angles d'un triangle est égale à $\pi$ donc
    \begin{align*}  
        &(\widehat{\overrightarrow{M A}, \overrightarrow{M\Omega}}) + (\widehat{\overrightarrow{A \Omega}, \overrightarrow{A M}}) + (\widehat{\overrightarrow{\Omega M}, \overrightarrow{\Omega A}}) \equiv \pi \mod 2\pi \\
        &\implies (\widehat{\overrightarrow{M A}, \overrightarrow{M\Omega}}) + (\widehat{\overrightarrow{A \Omega}, \overrightarrow{A M}}) \equiv (\widehat{\overrightarrow{\Omega A}, \overrightarrow{\Omega M}}) + \pi \mod 2\pi \\
        &\implies 2(\widehat{\overrightarrow{M A}, \overrightarrow{M\Omega}}) \equiv (\widehat{\overrightarrow{\Omega A}, \overrightarrow{\Omega M}}) + \pi \mod 2\pi
    \end{align*}
\end{demonstration}

\pagebreak

\begin{theoreme}{Théorème de l'angle inscrit au centre}{}
    Soit $\Gamma$ un cercle de centre $\Omega$.

    Pour tous $A, B, M \in \Gamma$ distincts, on a
    $$
    2(\widehat{\overrightarrow{MA}, \overrightarrow{MB}}) \equiv (\widehat{\overrightarrow{\Omega A}, \overrightarrow{\Omega B}}) \mod 2\pi
    $$
\end{theoreme}

\tikzexternalenable
\begin{figure}[H]
    \centering
    \resizebox{0.7\textwidth}{!}{%
    \begin{tikzpicture}[>=stealth, scale=1.3]

    % --- Cas 1 : Le centre est à l'intérieur de l'angle ---
    \begin{scope}[shift={(-2.5,0)}]
        \coordinate (O1) at (0,0);
        \draw[thick, gray] (O1) circle (2cm);
        \fill (O1) circle (1.5pt) node[below] {$\Omega$};

        \coordinate (A1) at (220:2cm);
        \coordinate (B1) at (320:2cm);
        \coordinate (M1) at (80:2cm);

        \fill (A1) circle (1.5pt) node[below left] {$A$};
        \fill (B1) circle (1.5pt) node[below right] {$B$};
        \fill (M1) circle (1.5pt) node[above] {$M$};

        \draw[->, thick, red!80!black] (M1) -- (A1);
        \draw[->, thick, red!80!black] (M1) -- (B1);
        \draw[->, thick, blue!80!black] (O1) -- (A1);
        \draw[->, thick, blue!80!black] (O1) -- (B1);

        \pic [draw, ->, "\small $\theta$", angle radius=0.9cm, angle eccentricity=1.3, red!80!black, thick] {angle = A1--M1--B1};
        \pic [draw, ->, "\small $2\theta$", angle radius=0.5cm, angle eccentricity=1.6, blue!80!black, thick] {angle = A1--O1--B1};
        
    \end{scope}

    % --- Cas 2 : Le centre est à l'extérieur de l'angle ---
    \begin{scope}[shift={(2.5,0)}]
        \coordinate (O2) at (0,0);
        \draw[thick, gray] (O2) circle (2cm);
        \fill (O2) circle (1.5pt) node[below left] {$\Omega$};

        % A et B du "même côté" pour que O soit à l'extérieur
        \coordinate (A2) at (10:2cm);
        \coordinate (B2) at (70:2cm);
        \coordinate (M2) at (160:2cm);

        \fill (A2) circle (1.5pt) node[right] {$A$};
        \fill (B2) circle (1.5pt) node[above right] {$B$};
        \fill (M2) circle (1.5pt) node[above left] {$M$};

        \draw[->, thick, red!80!black] (M2) -- (A2);
        \draw[->, thick, red!80!black] (M2) -- (B2);
        \draw[->, thick, blue!80!black] (O2) -- (A2);
        \draw[->, thick, blue!80!black] (O2) -- (B2);

        \pic [draw, ->, "\small $\theta$", angle radius=1.3cm, angle eccentricity=1.2, red!80!black, thick] {angle = A2--M2--B2};
        \pic [draw, ->, "\small $2\theta$", angle radius=0.6cm, angle eccentricity=1.4, blue!80!black, thick] {angle = A2--O2--B2};
    \end{scope}

\end{tikzpicture}
    }
\end{figure}
\tikzexternaldisable

\begin{demonstration}
    On a
    $$
    2(\widehat{\overrightarrow{MA}, \overrightarrow{MB}}) \equiv 2(\widehat{\overrightarrow{MA}, \overrightarrow{M\Omega}}) + 2(\widehat{\overrightarrow{M\Omega}, \overrightarrow{MB}}) \mod 2\pi
    $$

    Or $\Omega M A$ et $\Omega M B$ sont des triangles isocèles en $\Omega$ car $A$, $B$ et $M$ sont sur le cercle de centre $\Omega$.

    Ainsi, en appliquant le lemme précédent
    $$
    2(\widehat{\overrightarrow{MA}, \overrightarrow{M\Omega}}) \equiv (\widehat{\overrightarrow{\Omega A}, \overrightarrow{\Omega M}}) + \pi \mod 2\pi
    $$
    et
    $$
    2(\widehat{\overrightarrow{M\Omega}, \overrightarrow{MB}}) \equiv (\widehat{\overrightarrow{\Omega M}, \overrightarrow{\Omega B}}) + \pi \mod 2\pi
    $$

    Puis par Chasles
    \begin{align*}
        2(\widehat{\overrightarrow{MA}, \overrightarrow{MB}}) &\equiv (\widehat{\overrightarrow{\Omega A}, \overrightarrow{\Omega M}}) + (\widehat{\overrightarrow{\Omega M}, \overrightarrow{\Omega B}}) + 2\pi \mod 2\pi \\
        &\equiv (\widehat{\overrightarrow{\Omega A}, \overrightarrow{\Omega B}}) + 2\pi \mod 2\pi \\
        &\equiv (\widehat{\overrightarrow{\Omega A}, \overrightarrow{\Omega B}}) \mod 2\pi
    \end{align*}
\end{demonstration}

\begin{theoreme}{Condition de cocyclicité}{}
    Soient $A$, $B$, $C$ et $D$ quatre points, trois à trois non alignés, dans un plan vectoriel euclidien orienté.

    Ces quatre points sont cocycliques si et seulement si
    $$
    (\widehat{(CA), (CB)}) \equiv (\widehat{(DA), (DB)}) \mod \pi
    $$
\end{theoreme}

\tikzexternalenable
\begin{figure}[H]
    \centering
    \resizebox{0.9\textwidth}{!}{%
    \begin{tikzpicture}[>=stealth, scale=1.3]

    % --- Cas 1 : Quadrilatère convexe (Carré) ---
    \begin{scope}[shift={(-2.8,0)}]
        \coordinate (O1) at (0,0);
        \draw[dashed, gray] (O1) circle (2cm);

        % Points ordonnés sur le cercle
        \coordinate (A1) at (210:2cm);
        \coordinate (B1) at (330:2cm);
        \coordinate (C1) at (120:2cm);
        \coordinate (D1) at (50:2cm);

        \fill (A1) circle (1.5pt) node[below left] {$A$};
        \fill (B1) circle (1.5pt) node[below right] {$B$};
        \fill (C1) circle (1.5pt) node[above left] {$C$};
        \fill (D1) circle (1.5pt) node[above right] {$D$};

        % Droites prolongées avec calc
        \draw[thick, red!80!black] ($(A1)!-0.2!(C1)$) -- ($(C1)!-0.3!(A1)$);
        \draw[thick, red!80!black] ($(B1)!-0.2!(C1)$) -- ($(C1)!-0.3!(B1)$);

        \draw[thick, blue!80!black] ($(A1)!-0.2!(D1)$) -- ($(D1)!-0.3!(A1)$);
        \draw[thick, blue!80!black] ($(B1)!-0.2!(D1)$) -- ($(D1)!-0.3!(B1)$);

        % Angles
        \pic [draw, ->, angle radius=0.7cm, angle eccentricity=1.4, red!80!black, thick, pic text={\small $\alpha$}] {angle = A1--C1--B1};
        \pic [draw, ->, angle radius=0.7cm, angle eccentricity=1.4, blue!80!black, thick, pic text={\small $\alpha$}] {angle = A1--D1--B1};
        
    \end{scope}

    % --- Cas 2 : Quadrilatère croisé (Papillon) ---
    \begin{scope}[shift={(2.8,0)}]
        \coordinate (O2) at (0,0);
        \draw[dashed, gray] (O2) circle (2cm);

        % D est basculé de l'autre côté de la corde [AB]
        \coordinate (A2) at (210:2cm);
        \coordinate (B2) at (330:2cm);
        \coordinate (C2) at (120:2cm);
        \coordinate (D2) at (250:2cm);

        \fill (A2) circle (1.5pt) node[left] {$A$};
        \fill (B2) circle (1.5pt) node[right] {$B$};
        \fill (C2) circle (1.5pt) node[above left] {$C$};
        \fill (D2) circle (1.5pt) node[below] {$D$};

        % Droites prolongées
        \draw[thick, red!80!black] ($(A2)!-0.2!(C2)$) -- ($(C2)!-0.3!(A2)$);
        \draw[thick, red!80!black] ($(B2)!-0.2!(C2)$) -- ($(C2)!-0.3!(B2)$);

        \draw[thick, blue!80!black] ($(A2)!-0.2!(D2)$) -- ($(D2)!-0.3!(A2)$);
        \draw[thick, blue!80!black] ($(B2)!-0.2!(D2)$) -- ($(D2)!-0.3!(B2)$);

        % Angles
        \pic [draw, ->, angle radius=0.7cm, angle eccentricity=1.4, red!80!black, thick, pic text={\small $\alpha$}] {angle = A2--C2--B2};
        \pic [draw, <-, angle radius=0.4cm, angle eccentricity=1.6, blue!80!black, thick, pic text={\small $\alpha$}] {angle = B2--D2--A2};
        
    \end{scope}

\end{tikzpicture}
    }
\end{figure}
\tikzexternaldisable

\begin{demonstration}
    \begin{description}
        \item[Sens direct.]
        Soit $\Omega$ le centre du cercle passant par $A$, $B$, $C$ et $D$.

        Par le théorème de l'angle inscrit au centre appliqué $A$, $B$, $C$ puis à $A$, $B$, $D$, on trouve
        $$
        2(\widehat{\overrightarrow{CA}, \overrightarrow{CB}}) \equiv (\widehat{\overrightarrow{\Omega A}, \overrightarrow{\Omega B}}) \mod 2\pi
        $$
        et
        $$
        2(\widehat{\overrightarrow{DA}, \overrightarrow{DB}}) \equiv (\widehat{\overrightarrow{\Omega A}, \overrightarrow{\Omega B}}) \mod 2\pi
        $$
        Ainsi
        \begin{align*}
            &2(\widehat{\overrightarrow{CA}, \overrightarrow{CB}}) &\equiv 2(\widehat{\overrightarrow{DA}, \overrightarrow{DB}}) \mod 2\pi \\
            &\iff (\widehat{\overrightarrow{CA}, \overrightarrow{CB}}) &\equiv (\widehat{\overrightarrow{DA}, \overrightarrow{DB}}) \mod \pi \\
            &\iff (\widehat{(CA), (CB)}) &\equiv (\widehat{(DA), (DB)}) \mod \pi
        \end{align*}

        \item[Sens réciproque.]
        Supposons que $(\widehat{(CA), (CB)}) \equiv (\widehat{(DA), (DB)}) \mod \pi$.

        Il existe un unique cercle passant par $A$, $B$ et $C$, de centre $\Omega$.

        Soit $D'$ le point d'intersection de ce cercle avec la droite $(AD)$.

        Puisque $A$, $B$, $C$ et $D'$ sont cocycliques, on a
        \begin{align*}
            &(\widehat{(CA), (CB)}) &\equiv (\widehat{(D'A), (D'B)}) \mod \pi \\
            &\iff (\widehat{(D'A), (D'B)}) &\equiv (\widehat{(DA), (DB)}) \mod \pi \\
            &\iff (\widehat{(D'A), (D'B)}) &\equiv (\widehat{(D'A), (DB)}) \mod \pi \\
            &\iff (\widehat{(D'B), (DB)}) &\equiv 0 \mod \pi \\
        \end{align*}
        Donc $(D'B) = (DB)$ et $D \in (D'A)$ donc $D = D'$, ainsi $D$ est sur le cercle de centre $\Omega$ passant par $A$, $B$ et $C$.
    \end{description}
\end{demonstration}

\end{document}